\chapter{Filters, Matching Networks, and Transformers}
\label{ch:filters}
\section{Why Filters Matter in Amateur Radio}
Filters select desired signals and reject unwanted ones. They appear in receiver
front ends, IF stages, audio stages, transmitter outputs, anti-aliasing paths,
power supplies, and interference cures. With the s-plane chapter in hand, a filter
is simply a deliberate placement of poles (to shape the passband) and zeros (to
notch the stopband).

\section{Basic Filter Types}
\begin{description}
\item[Low-pass] Passes low frequencies and attenuates high frequencies.
\item[High-pass] Passes high frequencies and attenuates low frequencies.
\item[Band-pass] Passes a frequency band.
\item[Band-stop or notch] Rejects a frequency band.
\end{description}

\begin{figure}[htbp]
\centering
\begin{circuitikz}
\draw (0,0) node[left]{in} to[R,l=\(R\)] (3,0)
      to[C,l=\(C\)] (3,-2) node[ground]{};
\draw (3,0) to[short,-o] (4,0) node[right]{out};
\end{circuitikz}
\caption{First-order RC low-pass filter.}
\end{figure}

\section{Cutoff, Slope, and Order}
For a first-order RC filter, \(f_c=1/(2\pi RC)\). Past its corner, each pole
changes the asymptotic magnitude by about \(-\SI{20}{\decibel}\) per decade, and
each zero by \(+\SI{20}{\decibel}\) per decade. The \emph{order} of a filter is
its number of poles, so an \(n\)-pole low-pass rolls off at about
\(-20n\,\si{\decibel}\) per decade far above cutoff. Steeper skirts require more
poles.

\section{The Second-Order Section}
Higher-order filters are built from first-order and second-order sections. The
generic second-order low-pass transfer function is
\[
H(s)=\frac{\wzero^2}{s^2+\dfrac{\wzero}{Q}\,s+\wzero^2},
\]
identical in form to the series-RLC response of \cref{ch:splane}. Here the filter
\(Q\) is a design choice, not just a loss parameter: it sets how close the pole
pair sits to the \(\jj\omega\) axis and therefore how much the response peaks near
\(\wzero\). \(Q=1/\sqrt{2}\approx0.707\) gives the maximally flat (Butterworth)
second-order response with no peaking; higher \(Q\) peaks; lower \(Q\) rounds off
early. Band-pass and high-pass sections share the same denominator with different
numerators (a zero at the origin, or a double zero at the origin, respectively).

A passive second-order section needs both kinds of energy store---an inductor and
a capacitor (\cref{fig:lc-section}). The series \(L\) and shunt \(C\) give exactly
the \(H(s)\) above, with \(\wzero=1/\sqrt{LC}\) and a \(Q\) set by the source and
load resistances.

\begin{figure}[htbp]
\centering
\begin{circuitikz}
\draw (0,0) node[left]{in} to[L,l=\(L\)] (3,0)
      to[C,l=\(C\)] (3,-2) node[ground]{};
\draw (3,0) to[short,-o] (4,0) node[right]{out};
\end{circuitikz}
\caption{A passive second-order (LC) low-pass section: series inductor, shunt
capacitor, output across \(C\). Its transfer function is the \(H(s)\) above.}
\label{fig:lc-section}
\end{figure}

\section{Filter Families as Pole Patterns}
The named families differ only in \emph{where} the poles (and zeros) are placed
for a given order.

\begin{tabularx}{\textwidth}{@{}L{0.19\textwidth}Y@{}}
\toprule
Family & Pole/zero placement and consequence \\
\midrule
Butterworth & Poles equally spaced on a semicircle of radius \(\wzero\) in the left half-plane. Maximally flat passband, no ripple, moderate transition. \\
Chebyshev & Poles on an ellipse pushed closer to the \(\jj\omega\) axis. Sharper cutoff at the price of passband ripple and worse phase/transient behavior. \\
Elliptic (Cauer; ``elliptical'' in the pool) & Chebyshev-like poles plus transmission \emph{zeros} on the \(\jj\omega\) axis in the stopband. Sharpest transition; ripple in both bands. \\
Bessel & Poles positioned for maximally flat \emph{group delay} (linear phase); gentle amplitude cutoff. Good for pulse fidelity. \\
Crystal / cavity & Very high-\(Q\) resonators (poles hugging the axis) for narrow, highly selective band-pass shapes. \\
DSP FIR & No feedback, so no poles away from the origin (all zeros); can be exactly linear phase. \\
DSP IIR & Recursive: has poles, so it mimics analog families efficiently but with nonlinear phase. \\
\bottomrule
\end{tabularx}

\begin{figure}[htbp]
\centering
\begin{tikzpicture}[>=Latex,scale=1.15]
\draw[->,black!70] (-2.7,0) -- (0.8,0) node[right]{\(\sigma\)};
\draw[->,black!70] (0,-2.0) -- (0,2.0) node[above]{\(\jj\omega\)};
\draw[GuideLine,dashed] (0,0) circle (1.6);
% 4th-order Butterworth poles: angles 112.5,157.5,202.5,247.5 deg
\foreach \a in {112.5,157.5,202.5,247.5}{
  \filldraw[GuideBlue] ({1.6*cos(\a)},{1.6*sin(\a)}) circle (2.2pt);
}
\node[GuideLine] at (1.15,1.35) {\small \(|s|=\wzero\)};
\node[GuideBlue] at (-2.15,1.2) {\small Butterworth};
\end{tikzpicture}
\caption{Fourth-order Butterworth poles: four left-half-plane poles equally
spaced on the circle of radius \(\wzero\). Higher order adds more poles on the
same circle; Chebyshev pulls them toward the \(\jj\omega\) axis for a sharper but
rippled response.}
\label{fig:butterworth-poles}
\end{figure}

The \emph{transition band} is the frequency span between the passband edge and the
stopband; a ``moderate'' transition (Butterworth) is wider than the ``sharp''
transition of a Chebyshev or elliptic filter of the same order.
\Cref{fig:filter-families} shows all four analog families at the same order, which
makes the flatness-versus-sharpness-versus-ripple trade concrete.

\begin{figure}[htbp]
\centering
\includegraphics[width=0.9\textwidth]{filter_families}
\caption{Order-5 low-pass responses. Butterworth is maximally flat with a moderate
transition; Chebyshev trades passband ripple for a sharper skirt; elliptic is
sharpest of all, with ripple in both bands and stopband nulls (its
\(\jj\omega\)-axis zeros); Bessel is gentlest, chosen for clean phase and pulse
fidelity rather than sharpness.}
\label{fig:filter-families}
\end{figure}

\begin{controlsbox}
\textbf{Refresher: reading a filter as poles.} Filter design is pole placement in
reverse---you choose a target magnitude response and solve for pole locations.
Butterworth = poles on a circle (flattest); Chebyshev = poles on an ellipse
nearer the axis (higher per-stage \(Q\), sharper, rippled); elliptic adds
\(\jj\omega\)-axis zeros for stopband notches. The trade among flatness,
sharpness, and phase linearity is entirely a trade about how close the poles sit
to the imaginary axis.
\end{controlsbox}

\section{Impedance Matching}
An impedance-matching network does two jobs:
\begin{enumerate}
\item cancel the reactive part of a complex load, and
\item transform the remaining resistance to the desired system impedance.
\end{enumerate}
In a \(\SI{50}{\ohm}\) RF system, matching maximizes power transfer and keeps
transmission-line reflections controlled.

\subsection{The L-Network}
The simplest matching network uses two reactances. To match two real resistances
\(R_{\text{hi}}>R_{\text{lo}}\), place a series reactance on the low-resistance
side and a shunt reactance on the high-resistance side. The network operating
\(Q\) is set by the resistance ratio:
\[
Q=\sqrt{\frac{R_{\text{hi}}}{R_{\text{lo}}}-1},
\qquad
X_{\text{series}}=Q\,R_{\text{lo}},
\qquad
X_{\text{shunt}}=\frac{R_{\text{hi}}}{Q}.
\]
The two reactances are of opposite type (one inductor, one capacitor); choosing
which is the inductor sets whether the network is low-pass or high-pass
(\cref{fig:lnetwork}). An L-network has no free parameter left over, so its
bandwidth is fixed once the resistances are chosen.

\begin{figure}[htbp]
\centering
\begin{circuitikz}
\draw (0,0) node[left]{\(R_{\text{lo}}\)} to[L,l=\(X_{\text{series}}\)] (3.2,0);
\draw (3.2,0) to[C,l=\(X_{\text{shunt}}\)] (3.2,-2) node[ground]{};
\draw (3.2,0) to[short,-o] (4.4,0) node[right]{\(R_{\text{hi}}\)};
\end{circuitikz}
\caption{An L-network matching a low resistance to a high one: series reactance on
the low side, shunt reactance on the high side. Swapping which element is the
inductor makes it high-pass instead of low-pass.}
\label{fig:lnetwork}
\end{figure}

\subsection{Pi and T Networks}
Adding a third element gives a \(\pi\) or T network, which can be viewed as two
L-networks back to back that meet at an internal \emph{virtual resistance}
\(R_v\). Choosing \(R_v\) below both terminations sets a loaded \(Q\) of roughly
\[
Q_{\text{loaded}}\approx\sqrt{\frac{R_{\max}}{R_v}-1},
\]
so \(R_v\) is the extra degree of freedom that lets the designer pick the loaded
\(Q\)---and therefore the bandwidth and harmonic rejection---independently of the
resistance ratio. This is why transmitter output and antenna-tuner networks are
usually \(\pi\) or T: they match \emph{and} set selectivity. Higher loaded \(Q\)
narrows bandwidth and improves harmonic suppression at the cost of higher
circulating currents and losses.

\begin{mistakebox}
Matching is not the same thing as filtering. A network can do both, but the
concepts are distinct: matching controls the impedance seen by a source or line,
while filtering controls frequency-dependent transmission. A \(\pi\)-network tuner
happens to be a low-pass filter too, which is a feature, not the definition.
\end{mistakebox}

\section{Transformers}
For an ideal transformer,
\[
\frac{V_s}{V_p}=\frac{N_s}{N_p},
\qquad
\frac{I_s}{I_p}=\frac{N_p}{N_s},
\qquad
Z_p=\left(\frac{N_p}{N_s}\right)^2 Z_s.
\]
Impedance transforms as the turns ratio \emph{squared}. \textbf{Core saturation}
should be avoided: every magnetic core has a maximum flux density it can carry,
and once the drive pushes it near that limit the core's permeability collapses, so
the inductance drops and the winding current shoots up. The result is distortion,
heating, and loss of the intended impedance transformation---the reason
transformers and inductors are rated for a maximum voltage-time product (or power)
at a given frequency.

\section{Worked Example: Quarter-Wave and L-Network}
To match a real \(\SI{200}{\ohm}\) load to a \(\SI{50}{\ohm}\) system with a
lossless quarter-wave transformer,
\[
Z_t=\sqrt{Z_0Z_L}=\sqrt{(50)(200)}=\SI{100}{\ohm}.
\]
The same match with an L-network gives
\[
Q=\sqrt{\frac{200}{50}-1}=\sqrt{3}\approx1.73,
\qquad
X_{\text{series}}=1.73\times50\approx\SI{87}{\ohm},
\qquad
X_{\text{shunt}}=\frac{200}{1.73}\approx\SI{116}{\ohm}.
\]
Both realize the match; the quarter-wave version needs a line one quarter
wavelength long at the design frequency, while the L-network is lumped and
narrower band.
